Aborder la pratique historiquement informée de la musique - de la recherche de sources patrimoniales à la scène - implique d'analyser les besoins des chercheurs et musiciens. S'interroger sur leurs besoins permet de mettre en lumière le rôle prépondérant des bibliothèques musicales et d'orchestre dans la quête de la \enquote{résurrection de la musique du passé}\footnote{Landowska (Wanda), Musique ancienne, Paris, 1909}. (Wanda Landowska).

\subsection*{La pratique historiquement informée : des besoins particuliers pour les musiciens et les chercheurs.}

\begin{quote}
	De modestes tentatives d'une résurrection de la musique du passé ont été faites dans la deuxième moitié du dix-neuvième siècle, mais elles ont pris ces temps derniers des proportions admirables\footnote{Landowska (Wanda), Musique ancienne, Paris, 1909}. (Wanda Landowska)
\end{quote}

La \enquote{résurrection de la musique du passé} : voilà en quelques mots comment nous pourrions décrire l'esprit de la pratique historiquement informée de la musique. Ce courant est né dans les milieux académiques anglophones sous le nom de \textit{Historically Informed Practice} (HIP) suivant la proposition du critique musical Andrew Porter qui a été popularisée par l'ouvrage de John Butt, \textit{Playing with History : the historical approach to musical performance}\footnote{Butt (John), Playing with History, The historical approach to musical performance, Cambridge, Cambridge University Press, 2002.}. 
Les musiciens adeptes de cette pratique sont parfois appelés les \enquote{baroqueux} comme en témoigne le livre de Renaud Machart : \textit{Les baroqueux, un demi-siècle de musique, 1949-2001} \footnote{Machart (Renaud), Les baroqueux, un demi-siècle de musique, 1949-2001, Paris, Editions Fugue, 2024.}. 

La pratique historiquement informée est donc la pratique musicale qui consiste à nourrir son interprétation musicale en ayant une connaissance théorique et stylistique de l'époque du compositeur, à la fois musicale et matérielle. Cela consiste à jouer sur instruments d'époque, ou sur des copies d'instruments fidèles, en privilégiant les cordes en boyaux aux cordes synthétiques pour les instruments à cordes. Un soin tout particulier est donc apporté à l'utilisation des sources musicales patrimoniales telles que les partitions manuscrites autographes des compositeurs, ou encore des traités d'interprétation d'un instrument datant de la même époque ainsi qu'à des sources sur la facture instrumentale de l'époque.

Cette pratique a transformé la vie musicale du XX$^{e}$ siècle en contestant l'idée d'un progrès continu des formes et sonorités musicales :
\begin{quote}
	Jusqu'à présent, on n'a accordé que très peu d'attention aux transformations continuelles de la pratique musicale, les tenant même pour secondaires. La faute en incombe à cette conception qui voit un \enquote{développement} à partir des formes originales primitives, passant par des étapes intermédiaires plus ou moins déficientes, jusqu'à leur forme définitive \enquote{idéale}, laquelle est alors supérieure en tout point aux \enquote{étapes préliminaires}. [...] Mais depuis que nous sommes en mesure d'avoir une vue d'ensemble, cette opinion en ce qui concerne la musique, s'est trouvée réfutée : nous ne pouvons plus établir de différence de valeur entre la musique de Brahms, Mozart, Josquin ou Dufay - la théorie du progrès n'est plus défendable\footnote{Harnoncourt (Nikolaus),trad. Collin (Dennis), Le discours musical, pour une nouvelle conception de la musique, Paris, Gallimard, 1984. page 17.}. (Nikolaus Harnoncourt).
\end{quote}
Au contraire, la pratique historiquement informée analyse les différences des périodes pour en interpréter le plus fidèlement possible les œuvres. Cette pratique bouleverse donc la façon de travailler du musicien qui étudie les sources musicales patrimoniales pour nourrir son interprétation. 
La bibliothèque d'orchestre joue dès lors un rôle prépondérant dans les différentes phases d'une production musicale, depuis l'étude des sources patrimoniales, qui permet la compréhension la plus fine possible du contexte d'écriture, à la sélection des bonnes éditions du matériel (\enquote{les éditions Urtext tentent de restituer la version authentique, originale, des œuvres.}\footnote{Définition issue d'un article sur l'interprétation musicale et les éditions Urtext, en ligne à cette adresse : \url{https://lalettredumusicien.fr/article/urtext-et-interpretation-3188}}), qui accompagnent le musicien lors de la performance musicale historiquement informée. 

Comme l'a souligné Benoît Haug dans la postface du récent ouvrage \textit{Le son des musiques anciennes (1880-1980), imaginer, fabriquer, partager}\footnote{Gétreau (Florence),Framboisier (Alban), His (Isabelle), Le son des musiques anciennes (1880-1980), Imaginer, fabriquer, Rennes, Presses Universitaires de Rennes, 2024}, l'histoire de la pratique historiquement informée n'est pas linéaire mais se nourrit de plusieurs influences. 

Aux origines de la pratique historiquement informée, citons les pionniers de la fin du XIX$^{e}$ siècle et du début du XX$^{e}$ siècle Arnold Dolmetsch, Wanda Landowska et Geneviève Thibault de Chambure qui prônaient un retour aux instruments anciens comme le clavecin ou la viole de gambe dans un rejet de l'esthétique du romantisme tardif et un goût de l'exotisme. 
Si la génération actuelle des musiciens ne prétend pas en être héritière, elle se place plutôt dans la lignée des grands noms de la pratique historiquement informée du milieu du XXè siècle, tels Nikolaus Harnoncourt ou Gustav Leonhardt. Nés respectivement en 1929 et 1928, ces musiciens ont été témoins et acteurs du basculement de la pratique historiquement informée qui devient, entre les années 1950 et 1970, un projet artistique, scientifique et intellectuel à part entière. Leur période d'activité a d'ailleurs été propice à la création d'ensembles majeurs de la scène de la musique historique toujours en activité aujourd'hui, à l'image du Concentus Musicus Wien, créé par Nikolaus Harnoncourt, de l'ensemble Les Arts Florissants créé par William Christie ou encore de La Chapelle Royale créé par Philippe Herreweghe.

La diffusion de la pratique historiquement informée n'est pas simplement la transmission d'un état d'esprit, c'est le changement concret de la pratique de la musique ancienne. Elle s'est matérialisé notamment par l'ouverture de cursus universitaires comme le master d'interprétation des musiques anciennes de la Sorbonne (dont l'École des chartes est partenaire du parcours musique médiévale), ou la classe de pratique historiquement informée au CNSMDP \footnote{Conservatoire National Supérieur de Musique et Danse de Paris}, ou encore la formation de la Schola Cantorum de Bâle pour ne citer que les plus prestigieuses. Le travail sur les sources patrimoniales plébiscité par la pratique de la musique historique trouve un écho avec la démocratisation des travaux de recherche pour les élèves musiciens des grandes écoles dans le schéma LMD des études universitaires. Ainsi c'est tout l'univers musical qui est repensé pour la nouvelle génération de musiciens, dans les pas des précurseurs du XXè siècle. 


Il semble intéressant d'étudier les besoins particuliers des musiciens et des chercheurs adeptes de la pratique historiquement informée, dans leur quête du geste et de la technique musicale pour retrouver les sons du passé. 

L'article \enquote{L'interprétation de la musique historique}, issu de l'ouvrage de Nikolaus Harnoncourt \textit{Le Discours musical}\footnote{\textit{Ibid, pages 14 à 20.}} et écrit en 1954 est le premier écrit du chef sur le thème de la musique historique et le texte fondateur de son ensemble, le Concentus Musicus Wien. Il présente les enjeux des interprétations de la pratique historiquement informée. 
S'il y a deux façons de jouer de la musique historique, l'une en transposant la musique du passé au présent - en vigueur dans la musique jusqu'au XIX$^{e}$ siècle, l'autre en se replaçant dans l'époque d'écriture d'une œuvre musicale expérimentée dès le milieu du XXème siècle, c'est bien la deuxième option qui a été retenue par les pionniers de la musique historique et leurs héritiers. Leur attachement à la conformité de l'esprit de création et la restitution prétendument fidèle à l'œuvre sont les caractéristiques du jeu de ces personnalités musicales ayant fondé la discipline. 

\begin{quote}
	La volonté du compositeur est pour nous l'autorité suprême ; nous voyons la musique ancienne en tant que telle, dans sa propre époque et devons nous efforcer de la restituer authentiquement. [...] Mais une restitution est fidèle à partir du moment où elle s'approche de la conception qu'avait le compositeur au moment de la composition. On voit que ce n'est réalisable que jusqu'à un certain point : l'idée première d'une œuvre ne se laisse que deviner, en particulier lorsqu'il s'agit de musique d'époques très reculées. Les indices qui nous révèlent la volonté du compositeur sont les indictions d'exécution, l'instrumentation et les nombreux usages de la pratique d'exécution, en évolution constante, que le compositeur supposait naturellement connus de ses contemporains\footnote{Harnoncourt (Nikolaus), trad. Collin (Dennis), Le discours musical, pour une nouvelle conception de la musique, Paris, Gallimard, 1984. pages 16-17.}. (Nikolaus Harnoncourt)
\end{quote}

De ces mots de Nikolaus Harnoncourt découlent naturellement les besoins des chercheurs et des musiciens sensibles à la pratique historiquement informée. A l'instar de Gustav Leonhardt, qui passait de nombreuses heures en bibliothèque, les musiciens pratiquant la musique historique se doivent de préparer leurs interprétations en bibliothèques musicales. Les éléments importants auxquels les musiciens sont particulièrement attentifs sont, bien sûr, les \enquote{connaissances théoriques} des traités instrumentaux ou des traités de technique, \enquote{l'aspect formel et structurel de la restitution}, \enquote{l'image sonore (timbre, caractère, puissance des instruments)}, \enquote{la lecture de la notation et la pratique de l'improvisation} ainsi que le \enquote{concept et l'idéal sonore} voulu par le compositeur \footnote{Harnoncourt (Nikolaus), trad. Collin (Dennis), Le discours musical, pour une nouvelle conception de la musique, Paris, Gallimard, 1984. page 18.}.
\\\\
Ainsi, le premier besoin du musicien est d'accéder au texte original de l'œuvre musicale, c'est-à-dire, autant que possible, au manuscrit autographe.

\begin{quote}
	La notation tant des différentes parties que, tout particulièrement, de la partition, dégage outre le contenu purement informatif, un rayonnement suggestif, une magie à laquelle aucun musicien sensible n peut échapper [...]. Bien que ce rayonnement émane aussi des textes musicaux imprimés, il est manifestement le propre, dans une beaucoup plus grande mesure, des musiques manuscrites, étant bien entendu le plus fort lorsqu'il s'agit d'un autographe, du manuscrit du compositeur. [...] C'est pourquoi il est pour nous de la plus grande importance de découvrir l'œuvre que nous exécutons dans sa forme originale, c'est-à-dire de recourir autant que possible à la partition originale ou du moins à une bonne édition en fac-similé \footnote{Harnoncourt (Nikolaus), trad. Collin (Dennis), Ce que dit un autographe in Le discours musical, pour une nouvelle conception de la musique, Paris, Gallimard, 1984. (page 244)}.(Nikolaus Harnoncourt)
\end{quote}

L'édition originale d'une œuvre et davantage son manuscrit permettent au musicien de découvrir la méthode de travail du compositeur et donne également des clés de compréhension de son œuvre. Elles lui sont précieuses pour comprendre les variations dans la notation musicale, discipline qui a tant changé qu'il est parfois compliqué de lire la musique d'époque sans avoir de notions de paléographie musicale. 
Le musicien peut trouver ce type de sources dans les départements de la musique des bibliothèques nationales, qui sont des bibliothèques musicales à part entière, de préférence sur leurs outils en ligne (comme Gallica pour la Bibliothèque nationale de France).
Au-delà des partitions originales, les musiciens sont également à la recherche des traités d'époque pour apprendre la technique d'une période et d'un instrument donnés. Cela leur permet de perfectionner leur jeu dans l'esprit de l'œuvre qu'ils souhaitent interpréter. Les traités des XVIIè et XVIIIè siècles, la méthode de violon de Léopold Mozart par exemple, le traité de flûte de Quantz, ou encore les traités de Rameau et de C.P.E Bach, sont particulièrement riches pour les musiciens. 

Si les traités permettent aux musiciens d'apprendre la technique de jeu d'une période donnée, le type d'instrument utilisé est également très important pour parvenir à restituer l'esprit d'une œuvre ou de son époque. Il est impossible de jouer une difficulté d'une époque avec un instrument qui n'y est pas adapté. C'est pour cela que la pratique historiquement informée s'accompagne d'un courant de facture instrumentale : \enquote{L'archéologie musicale} dont l'objet est de connaître les instruments anciens afin de pouvoir les restaurer ou les imiter dans un objectif d'authenticité de la restitution musicale. Le travail de \enquote{l'archéologie musicale}, dont les précurseurs  Auguste Tolbecque, Arnold Dolmetsch, Théodore Reinach et Victor-Charles Mahillon ont dressé les bases de la méthodologie, se fonde sur trois sources : l'archéologie, l'étude des textes et l'étude de l'iconographie \footnote{Rault (Christian), Auguste Tolbecque et l'archéologie musicale expérimentale : des institutions d'un pionnier aux pratiques contemporaines, in Le son des musiques anciennes (1880-1980), Imaginer, fabriquer, partager, Rennes, Presses Universitaires de Rennes, 2024. (pages 245 à 256).}. Cette exploitation des sources se manifeste donc par un important travail de recherche en musées et en collections privées sur la facture instrumentale, et c'est ainsi que les musiciens adeptes de la pratique historiquement informée, guidés par leurs trouvailles en bibliothèques, \enquote{prônent [...] la recherche des sonorités authentiques, imposent l'usage d'instruments d'époque, avec cordes en boyau, des diapasons différents}\footnote{Rault (Christian), Auguste Tolbecque et l'archéologie musicale expérimentale : des institutions d'un pionnier aux pratiques contemporaines, dans Le son des musiques anciennes (1880-1980), Imaginer, fabriquer, partager, Rennes, Presses Universitaires de Rennes, 2024. (page 246)}.

Il faut alors remarquer que le travail des musiciens adeptes de la pratique historiquement informée est très lié au travail des musicologues du même courant. 
Les chercheurs spécialistes, d'une époque ou d'une oeuvre, sont en effet des personnes-ressources pour les musiciens en quête d'authenticité, lorsqu'ils travaillent leur interprétation.
En effet,  les chercheurs spécialisés en organologie, iconographie musicale, ou tout simplement spécialistes d'une période ou d'un compositeur, peuvent prodiguer de précieux conseils à un musicien ou à un chef qui chercherait à construire une interprétation.
Leurs compétences très spécialisées leur permettent par exemple de faire le rapprochement entre la facture instrumentale d'un violon exposé dans un musée de la musique, et le violon pour lequel une œuvre a été dédiée. Les deux éléments sont précisément expliqués dans le manuscrit autographe original de l'œuvre qu'un musicien souhaiterait faire reproduire afin d'interpréter l'œuvre. 

Dans cette conception qui préconise la construction méticuleuse de l'interprétation, la bibliothèque d'orchestre se voit attribuer un nouveau rôle qui dépasse la préparation intellectuelle et physique d'une production (choix éditorial, location ou achat, coups d'archets, coupures...) puisqu'elle devient le centre névralgique de l'interprétation de la musique historiquement informée. 


\subsection*{Les bibliothèques musicales et d'orchestre, des sources patrimoniales utiles aux musiciens et chercheurs de la musique historique.}

Dans le contexte de développement de la pratique historiquement informée, le rôle des bibliothèques, musicales d'une part et d'orchestre d'autre part, revêt une importance nouvelle. 
S'éloignant progressivement du dépôt opérationnel d'archives où le bibliothécaire était là pour assurer la préparation, le suivi et la manutention des partitions durant une production, la bibliothèque devient le coeur de la musique historique. 

Ainsi, la bibliothèque musicale, à l'instar du département de la Musique de la Bibliothèque nationale de France, est le lieu de consultation de sources patrimoniales (sur site ou en ligne) privilégié des chercheurs.

Quant aux musiciens sensibles à la pratique historiquement informée, ils préparent leurs interprétations en héritiers des pionniers du travail en bibliothèque, Henri Casadesus, Louis Diémer et Geneviève Thibault de Chambure\footnote{Lafitte, (Priscille), animatrice. Geneviève de Chambure, pionnière du baroque.[Série de podcats audio] Les Sagas musicales, Radio France, 2026.\url{https://www.radiofrance.fr/francemusique/podcasts/serie-genevieve-de-chambure-pionniere-du-baroque}} qui étudiaient en bibliothèque des partitions des XVIIè et XVIIIè siècles afin de les interpréter sur instruments d'époque\footnote{Becker-Derex (Christiane), \enquote{La création des sociétés des instruments anciens : vers une nouvelle esthétique sonore}. Dans Le son des musiques anciennes (1880-1980), Imaginer, fabriquer, partager, éd. Florence Gétreau, Alban Framboisier et Isabelle His, Rennes, Presses universitaires de Rennes, 2024.}. Ils ne cherchent plus nécessairement à y trouver des partitions oubliées mais à y lire les annotations du compositeur ou de chefs pouvant les aider à se rapprocher au maximum du son authentique.

Ainsi, des séances d'étude en bibliothèque musicale, pour y étudier le manuscrit original, l'édition critique de musicologues spécialistes de pratique historiquement informée, ou encore les traités d'époque, permettent de réaliser l'interprétation la plus fidèle possible en prêtant une attention toute particulière aux éléments exigeant la vigilance de l'interprète. 
Dans son ouvrage \textit{Le Discours musical}\footnote{Harnoncourt (Nikolaus), trad. Collin (Dennis), Le discours musical, pour une nouvelle conception de la musique, Paris, Gallimard, 1984.}, Nikolaus Harnoncourt consacre un chapitre à chaque domaine auquel l'interprète doit être vigilant. Il s'agit de  \enquote{la notation\footnote{Ibid. Pages 34 à 49}}, \enquote{l'articulation\footnote{Ibid. Pages 50 à 65}}, \enquote{le mouvement\footnote{Ibid. Pages 66 à 77}} (c'est-à-dire le rythme et le tempo donnés à l'œuvre), ainsi que \enquote{le système sonore et l'intonation\footnote{Ibid. Pages 78 à 88}}.  

Ces enjeux de pratique historiquement informée soulignent le rôle fondamental des bibliothèques musicales dans la recherche et l'interprétation musicales. En effet, les fonds musicaux s'ils sont d'un \enquote{intérêt majeur pour l'historien\footnote{Hausfater (Dominique) et Campos (Rémy), \textit{L’interprétation musicale dans les fonds des bibliothèques}: Journée d’étude Du 7 Janvier 2006 Conservatoire de Musique de Genève / Haute École de Musique.” Fontes Artis Musicae, vol. 54, no. 1, 2007, pages. 1–5. JSTOR,\url{http://www.jstor.org/stable/23510565.} (consulté le 25 août 2026)}}, soulèvent également des questionnements bibliothéconomiques et musicologiques au niveau des recherches sur l'interprétation.





\subsection*{Étude de cas : LaDocumenta.eu, un corpus de référence pour les chercheurs et musiciens de la pratique historiquement informée.}

LaDocumenta.eu, en tant que bibliothèque numérique spécialisée en musique, est un corpus de référence pour les chercheurs et les musiciens de la pratique historiquement informée. Celle-ci rassemble les ressources numériques et numérisées des cinq orchestres européens membres du consortium LaDocumenta.

La plateforme a pour ambition de documenter la pratique historiquement informée de la musique en suivant toutes les phases d'une production musicale que sont la pré-production, la production, le concert et la post-production.

Elle dispose pour cela d'un corpus de 4835 ressources\footnote{Nombre de ressources au 13 août 2026} structuré autour de quarante-huit typologies de ressources.
\begin{landscape}
	\begin{table}[htbp]
		\centering
		\caption{Liste des 48 ressources classées en fonction de la phase de production.}
		\label{tab:ressources_production}
		
		\vspace{0.3cm}
		
		{\tiny
			\renewcommand{\arraystretch}{1.15}
			\setlength{\tabcolsep}{5pt}
			\setlength{\extrarowheight}{3pt}
			
			\begin{tabular}{|p{3.4cm}|p{6.8cm}|p{3.8cm}|p{3.2cm}|p{3.2cm}|}
				\hline
				\textbf{Pré-production (1)} \newline \textit{(production, préparation)} & 
				\textbf{Pré-production (2)} \newline \textit{(éducation, information)} & 
				\textbf{En-production} & 
				\textbf{Concerts} & 
				\textbf{Post-production} \\ 
				\hline
				
				• Déclaration d’intention \newline
				• Plan structurel du conducteur \newline
				• Concept structurel de la production \newline
				• Plan de répétition \newline
				• Document préparatoire divers \newline
				• Contenu dramaturgique \newline
				• Nomenclature & 
				
				• Conducteur original (source prim.) \newline
				• Partition originale (source prim.) \newline
				• Édition originale (source prim.) \newline
				• Livret original (source prim.) \newline
				• Nouvelle édition basée sur sources prim. (source sec.) \newline
				• Nouvelle partition basée sur nouvelle éd. (source sec.) \newline
				• Pitch original (source prim.) \newline
				• Tempérament original (source prim.) \newline
				• Illustration d’instrument (source prim./iconog.) \newline
				• Image d’instrument (photo/vidéo) (source sec./iconog.) \newline
				• Illustration d’ensemble, continuo, chef, position/posture (source prim.) \newline
				• Image d’ensemble, continuo, chef, position/posture (photo/vidéo) (source sec.) \newline
				• Extrait de traité instrumental (source prim.) \newline
				• Extrait de traité sur styles (source prim.) \newline
				• Extrait de livre musique hist. formée (source sec.) \newline
				• Image d’instrument de musée (source prim.) \newline
				• Style orig. d’interface utilisateur (archets, embouchures...) (source prim.) \newline
				• Instrument moderne ou interface conçue musique hist. formée (source sec.) \newline
				• Enregistrement (source prim.) & 
				
				• Comm. sur répertoire (rareté, pépite, fake news...) \newline
				• Plan de salle \newline
				• Ajustement (plan répétition, coupes, réorg., durée...) \newline
				• Compromis (instruments, pitchs, styles) \newline
				• Annotations de conducteur \newline
				• Annotations de partition \newline
				• Enregistrement de test \newline
				• Notes de programmes \newline
				• Photo répétitions (pupitres, interactions...) \newline
				• Interviews musiciens (audio, vidéo) \newline
				• Unpublished artist statement \newline
				• Vidéo matériel éducatif & 
				
				• Enregistrement final \newline
				• Photo concert / enregistrement public \newline
				• Matériel sortant (synthèse, kit presse...) \newline
				• Matériel éducatif de communication & 
				
				• Article de presse \newline
				• Partition avec annotations finales \newline
				• Bilan de production \newline
				• Évaluation plan de production \newline
				• Critique \newline
				• Conducteur avec annotations finales \\ 
				\hline
			\end{tabular}
		}
	\end{table}
\end{landscape}
Les phases du cycle de production musicale sont réparties, dans la plateforme, en quatre catégories que sont la pré-production, la production, le concert et la post-production.
La phase de pré-production répartit les ressources entre les thématiques d’éducation/information (qui rassemble les partitions, livrets et traités originaux ainsi que les éditions critiques et l’organologie) et la préparation d’une production (cheminement artistique en amont des répétitions, comme les nomenclatures…).
La phase de production, quant-à-elle, documente le travail des répétitions. Un soin particulier est donné aux plans de scène et annotations de travail sur les partitions par exemple.
Moment tant attendu du cycle de la production musicale, le concert est l’aboutissement du cycle de production. La plateforme LaDocumenta.eu permet la conservation des enregistrements officiels, des photographies ainsi que des dossiers de valorisation. 
Enfin, la phase de post-production conserve la mémoire technique et critique de la production, entre réception critique et retour de l’institution. Elle archive les partitions avec les annotations finales. Il existe en outre deux types de sources pour le corpus : les sources primaires, qui sont des matériaux historiques d’époque, comme des partitions annotées par les chefs, et les sources secondaires qui sont des outils d’analyse ou encore des restitutions. 




